\section{Discussion}
\label{sec:discussion}

\subsection{What the Numbers Do and Do Not Claim}

Near-perfect stage-1 performance (F1 $\approx 0.9995$) should not be mistaken for a claim that anomaly detection in a real plant is solved. \nppad{} is generated by PCTRAN, a deterministic simulator with idealized sensors and no noise model. Real plants exhibit sensor drift, instrument channels that fail independently of the process, off-nominal transients that are not accidents (grid frequency excursions, operator test actions, scheduled maintenance evolutions), and seasonal variations in heat-sink temperatures. Each of these routinely produces excursions that a naive reconstruction-error detector would flag as anomalous. The correct reading of the stage-1 result is that on clean simulator data, screening is a near-trivial problem and the informative difficulty lies further downstream. A production deployment of stage 1 would require noise-aware training data, a domain-adaptation step onto plant sensor streams, and calibration of false-positive cost against the cost of a missed early warning.

The $71.2\%$ stage-2 accuracy is more informative about what is achievable today. The failure modes split cleanly into two groups. The first is physical ambiguity on short windows: LOCA and LOCAC, for example, share nearly identical primary-side signatures in the first 30 seconds of the transient and diverge only once containment isolation logic activates. Extending the window, cascading across windows, or explicitly modeling containment-state features would likely close this gap. The second group is class imbalance: LACP, ATWS, and SP each have fewer than $100$ training windows, and no amount of architectural choice compensates for data scarcity of this kind. Data augmentation, severity-conditioned sampling, and class-balanced loss functions are natural next steps.

The $68.6\%$ reduction in point-kinetics residual and $39.6\%$ reduction in conservation violation from the physics-informed surrogate are robust because the comparison is an architecture-matched ablation---the two models are identical except for the physics-loss terms. The raw trajectory MSE is comparable between the two models; this is the expected behavior of a physics-informed regularizer. The practical value is not MSE reduction but the suppression of the heavy tail of physically implausible predictions, which is what makes the physics residual a useful consistency signal for the integrated pipeline.

\subsection{Limitations}

\paragraph{Simulator-to-plant gap.} All evaluations are on PCTRAN data. The simulator faithfully captures first-order neutron kinetics and thermal hydraulics but simplifies containment, secondary-side heat transfer, and control-system behavior relative to a real plant. Transfer to real plant data would require retraining with noise- and drift-augmented streams and would likely reveal a gap in stage-2 performance that is currently masked by the clean simulator distribution.

\paragraph{Simplified physics in stage 3.} The surrogate embeds point kinetics and steady-state energy conservation, not full spatial neutron transport, not two-phase thermal hydraulics with flow-regime-dependent closures, and not primary-to-secondary heat exchange with dynamic U-tube inventory. These simplifications are intentional for a plant-level, real-time-capable surrogate, but they mean the physics residual is a lower bound on physical inconsistency rather than a full-system residual. Accidents that violate the simplified physics will be flagged; accidents that violate only higher-order physics may slip through.

\paragraph{Cross-validation confounds surrogate uncertainty with classifier error.} Case study (iii) in Figure~\ref{fig:case-studies} shows that rare-class transients can produce high $\kappa$ because the surrogate itself is uncertain, not because the classifier is wrong. A principled fix is to attach uncertainty estimates (e.g., deep ensembles or MC dropout) to the surrogate and to interpret $\kappa$ only when surrogate uncertainty is below a threshold. We leave this to future work.

\paragraph{Window size and latency.} A 30-step window is adequate for classifying accidents whose signatures establish quickly but is a hard ceiling on how much history the classifier can exploit. For slowly developing events and for distinguishing physically similar accidents, a longer context combined with streaming inference would be preferable.

\subsection{Generalization to Advanced Reactor Designs}

The framework is designed around a small-molecular-reactor-friendly interface: a plant-level sensor window in, a screening flag and classification and physics-checked forecast out. Porting to a different reactor design requires changing the physics of stage 3 but not the shape of the pipeline. For sodium fast reactors (SFRs, e.g., TerraPower), the point-kinetics form survives but the reactivity coefficients change sign structure and magnitude, and the energy-conservation constraint must reflect sodium thermal properties. For molten salt reactors (MSRs, e.g., Kairos), delayed neutron precursors are transported by the flowing fuel and the point-kinetics equations gain an advective term that is central, not peripheral, to the physics loss. For small PWRs (Oklo Aurora-class PWR variants, NuScale integral PWRs), the equations of Section~\ref{sec:pinn} apply directly with updated physical parameters, and the main change is that SMR designs typically use passive safety systems whose signatures are under-represented in \nppad{}.

In all three cases, the diagnostic pipeline structure---screening, classification with attribution, physics-checked prognosis---is reactor-agnostic. Only the physics terms and the class taxonomy are reactor-specific.

\subsection{Toward Deployment}

A production system along these lines would require three additions beyond what is reported here. First, uncertainty quantification on both the classifier and the surrogate, so that low-confidence diagnoses can be surfaced as such rather than silently acted upon. Second, continuous domain adaptation between the simulator-trained models and the live plant data, likely via a small amount of on-plant fine-tuning plus distribution-shift monitoring. Third, human-factors integration: the SHAP attributions and physics-consistency scores are intended to support operator decision-making, and their presentation format must be co-designed with control-room personnel rather than delivered as raw heatmaps. None of these are research open problems; all are engineering problems with established solution paths.
